\labsection{10}{Reinforcement learning and final integration}{sec:lab10-reinforcement}

\begin{labproject}{Learn a policy by interaction, then connect the result to the full course workflow}
\textbf{Coverage.} Chapter 20 with a handoff to the capstone.\\
\textbf{Source files.} \texttt{Lab10/Task01.ipynb}.\\
\textbf{Core environment.} A self-contained $4\times4$ GridWorld.\\
\textbf{Goal.} Train a tabular Q-learning agent, evaluate its greedy policy separately from exploratory training, and articulate how reinforcement learning changes the data/label/evaluation assumptions used earlier in the course.

\begin{labmode}
This is the tenth and final scheduled lab. The capstone that follows is an integration project, not Lab 11.
\end{labmode}

\medskip\noindent\textbf{Part A --- Tabular Q-learning in GridWorld.}\par
\begin{lstlisting}[style=mlpython]
import numpy as np
import matplotlib.pyplot as plt
\end{lstlisting}

\begin{lstlisting}[style=mlpython]
rng = np.random.default_rng(42)

# A tiny 4x4 GridWorld. Start at the top-left, reach the goal at bottom-right,
# and avoid the trap. No external environment package is required.
ROWS, COLS = 4, 4
START = (0, 0)
GOAL = (3, 3)
TRAP = (1, 3)
ACTIONS = [(-1, 0), (1, 0), (0, -1), (0, 1)]  # up, down, left, right
ACTION_NAMES = ["up", "down", "left", "right"]

Q = np.zeros((ROWS, COLS, len(ACTIONS)), dtype=float)
alpha = 0.15
gamma = 0.95
epsilon = 1.0
min_epsilon = 0.05
\end{lstlisting}

\begin{lstlisting}[style=mlpython]
def step(state, action):
    r, c = state
    dr, dc = ACTIONS[action]
    nr = min(max(r + dr, 0), ROWS - 1)
    nc = min(max(c + dc, 0), COLS - 1)
    next_state = (nr, nc)

    if next_state == GOAL:
        return next_state, 1.0, True
    if next_state == TRAP:
        return next_state, -1.0, True
    return next_state, -0.04, False
\end{lstlisting}

\begin{lstlisting}[style=mlpython]
def choose_action(state, explore=True):
    if explore and rng.random() < epsilon:
        return int(rng.integers(len(ACTIONS)))
    r, c = state
    return int(np.argmax(Q[r, c]))
\end{lstlisting}

\begin{lstlisting}[style=mlpython]
training_returns = []

for episode in range(5000):
    state = START
    done = False
    episode_return = 0.0

    for _ in range(100):
        action = choose_action(state, explore=True)
        next_state, reward, done = step(state, action)
        episode_return += reward

        r, c = state
        nr, nc = next_state
        backup_value = reward if done else reward + gamma * np.max(Q[nr, nc])
        Q[r, c, action] += alpha * (backup_value - Q[r, c, action])
        state = next_state

        if done:
            break

    training_returns.append(episode_return)
    epsilon = max(min_epsilon, epsilon * 0.9985)
\end{lstlisting}

\begin{lstlisting}[style=mlpython]
window = 100
moving = np.convolve(training_returns, np.ones(window) / window, mode="valid")
plt.figure(figsize=(6, 4))
plt.plot(np.arange(window - 1, len(training_returns)), moving)
plt.xlabel("Episode")
plt.ylabel("Mean return over last 100 episodes")
plt.title("Q-learning training curve")
plt.tight_layout()
plt.show()
\end{lstlisting}

\begin{lstlisting}[style=mlpython]
# Evaluate the learned greedy policy.
successes = 0
returns = []
for _ in range(200):
    state = START
    total_reward = 0.0
    for _ in range(100):
        action = choose_action(state, explore=False)
        state, reward, done = step(state, action)
        total_reward += reward
        if done:
            successes += int(state == GOAL)
            break
    returns.append(total_reward)

print("success rate:", successes / 200)
print("mean return:", round(float(np.mean(returns)), 3))
print("learned action at start:", ACTION_NAMES[choose_action(START, explore=False)])
\end{lstlisting}

\begin{tryit}
Change one of reward, trap location, exploration schedule, or discount factor. Predict the behavioral consequence before retraining and compare the learned greedy policy afterward.
\end{tryit}
\end{labproject}

\begin{labdeliverable}
Submit the reward specification, one pre-training prediction, the learning curve, repeated greedy-policy success rate, and one controlled environment or hyperparameter change. End with a short capstone handoff naming the learning signal, validation boundary, baseline, and metric you would choose for a new problem.
\end{labdeliverable}
